% 第5章图：分叉树上的消息与回溯。树边{1,3},{2,3},{3,4},{4,5}，根3。
% phi1=(0,2),phi2=(2,0),phi3=(1,0),phi4=(0,1),phi5=(3,0)；不同取值成本1,1,2,1。
% m1->3=(0,1), m2->3=(1,0), m5->4=(1,0), m4->3=(1,1)；根(3,2)；回溯 x=(0,1,1,1,1)，成本2。
\begin{tikzpicture}[x=1mm,y=1mm,>=Stealth,
  v/.style={circle,draw=FigureStroke,fill=FigureBlueFill,line width=.8pt,minimum size=6.5mm,inner sep=0pt},
  w/.style={circle,draw=FigureStroke,fill=white,line width=.8pt,minimum size=6.5mm,inner sep=0pt},
  ord/.style={circle,draw=FigureMuted,line width=.5pt,inner sep=.4pt,minimum size=3.6mm,font=\FigureNoteFont\fontsize{7}{8}\selectfont,text=FigureMuted},
  msg/.style={font=\figurefont\footnotesize,text=FigureInk},
  note/.style={font=\FigureNoteFont\fontsize{7.5}{9}\selectfont,text=FigureMuted}]
% ---------- (a) upward messages ----------
\begin{scope}
 \node[v] (a3) at (18,26) {$3$};
 \node[v] (a1) at (0,8) {$1$};
 \node[v] (a2) at (18,8) {$2$};
 \node[v] (a4) at (36,8) {$4$};
 \node[v] (a5) at (36,-10) {$5$};
 \foreach \u/\v in {a1/a3,a2/a3,a4/a3,a5/a4}{\draw[->,FigureBlue,line width=.9pt] (\u)--(\v);}
 \node[msg,anchor=south east] at (8,18) {$(0,1)$};
 \node[ord] at (-2.5,20) {1};
 \node[msg,anchor=west] at (19.5,13) {$(1,0)$};
 \node[ord] at (15.4,13) {2};
 \node[msg,anchor=south west] at (28,18) {$(1,1)$};
 \node[ord] at (38.8,20) {4};
 \node[msg,anchor=west] at (37.5,-1) {$(1,0)$};
 \node[ord] at (33.4,-1) {3};
 \node[note,anchor=south] at (a3.north) {根：$(3,2)$};
 \node[note,anchor=north] at (a1.south) {$(0,2)$};
 \node[note,anchor=north] at (a2.south) {$(2,0)$};
 \node[note,anchor=west] at (a4.east) {$(0,1)$};
 \node[note,anchor=west] at (a5.east) {$(3,0)$};
 \node[note,anchor=east] at (a3.west) {$(1,0)$};
 \node[font=\figurefont\footnotesize] at (18,-20) {(a) 自叶向根的消息};
\end{scope}
% ---------- (b) backtracking ----------
\begin{scope}[xshift=54mm]
 \node[v] (b3) at (18,26) {$3$};
 \node[w] (b1) at (0,8) {$1$};
 \node[v] (b2) at (18,8) {$2$};
 \node[v] (b4) at (36,8) {$4$};
 \node[v] (b5) at (36,-10) {$5$};
 \foreach \u/\v in {b3/b1,b3/b2,b3/b4,b4/b5}{\draw[->,FigureYellow,line width=1pt] (\u)--(\v);}
 \node[msg,anchor=south east] at (8,18) {$1\mapsto0$};
 \node[msg,anchor=west] at (19.5,13) {$1\mapsto1$};
 \node[msg,anchor=south west] at (28,18) {$1\mapsto1$};
 \node[msg,anchor=west] at (37.5,-1) {$1\mapsto1$};
 \node[note,anchor=south] at (b3.north) {根取$x_3=1$};
 \node[note,anchor=north] at (b1.south) {$x_1=0$};
 \node[note,anchor=north] at (b2.south) {$x_2=1$};
 \node[note,anchor=west] at (b4.east) {$x_4=1$};
 \node[note,anchor=west] at (b5.east) {$x_5=1$};
 \node[font=\figurefont\footnotesize] at (18,-20) {(b) 自根向外的回溯};
\end{scope}
\end{tikzpicture}
